\documentclass{article}

\usepackage{amsmath,amssymb}
\usepackage{xurl}
\usepackage[hidelinks]{hyperref}

% Shared commands for notes in docs/paper-gaps/.

\DeclareUnicodeCharacter{2080}{\ifmmode{}_{0}\else\textsubscript{0}\fi}
\DeclareUnicodeCharacter{2081}{\ifmmode{}_{1}\else\textsubscript{1}\fi}
\DeclareUnicodeCharacter{2082}{\ifmmode{}_{2}\else\textsubscript{2}\fi}
\DeclareUnicodeCharacter{2099}{\ifmmode{}_{n}\else\textsubscript{n}\fi}

\newcommand{\Lean}{\textsc{Lean}}
\ExplSyntaxOn
\NewDocumentCommand{\leanid}{m}{
  \ifmmode
    \textsf{\detokenize{#1}}
  \else
    \group_begin:
    \urlstyle{sf}
    \tl_set:Nn \l_tmpa_tl {#1}
    \tl_replace_all:Nnn \l_tmpa_tl {\_} {_}
    \exp_args:NV \path \l_tmpa_tl
    \group_end:
  \fi
}
\ExplSyntaxOff
\providecommand{\tr}{\operatorname{tr}}
\newcommand{\tnleanIssueBase}{https://github.com/LionSR/TNLean/issues/}
\newcommand{\tnleanPrBase}{https://github.com/LionSR/TNLean/pull/}
\newcommand{\tnleanDocBase}{https://sirui-lu.com/TNLean/docs/}
\newcommand{\ghissue}[1]{\href{\tnleanIssueBase#1}{GitHub issue \##1}}
\newcommand{\ghpr}[1]{\href{\tnleanPrBase#1}{PR \##1}}
\newcommand{\doclink}[1]{\href{\tnleanDocBase#1.html}{\path{#1}}}

% Dirac notation and the identity operator, as in blueprint/src/macros/common.tex.
% \providecommand keeps note-local definitions authoritative.
\usepackage{dsfont}
\providecommand{\ket}[1]{|#1\rangle}
\providecommand{\bra}[1]{\langle #1|}
\providecommand{\braket}[2]{\langle #1 | #2 \rangle}
\providecommand{\ketbra}[2]{|#1\rangle\!\langle #2|}
\providecommand{\Id}{\mathds{1}}
\providecommand{\one}{\mathds{1}}
\providecommand{\C}{\mathbb{C}}
\providecommand{\MN}[1]{M_{#1}(\C)}
\providecommand{\MD}{M_D(\C)}

% External referencing for paper sources.  The build helper writes sanitized
% auxiliary files under build/paper-aux/ and excludes duplicated source labels.
\usepackage{xr}

% Import a generated paper auxiliary file with a caller-supplied label prefix.
% The candidate paths support builds from docs/paper-gaps/, the repository root,
% and build/paper-aux/ itself.
\newcommand{\importpaperaux}[2]{%
  \IfFileExists{../../build/paper-aux/#2.aux}{%
    \externaldocument[#1]{../../build/paper-aux/#2}%
  }{%
    \IfFileExists{build/paper-aux/#2.aux}{%
      \externaldocument[#1]{build/paper-aux/#2}%
    }{%
      \IfFileExists{#2.aux}{%
        \externaldocument[#1]{#2}%
      }{}%
    }%
  }%
}

% General external reference macros
\newcommand{\paperref}[2]{\ref{#1-#2}}
\newcommand{\papereqref}[2]{(\ref{#1-#2})}

% CPSV16 source (1606.00608 / MPDO-22-12-17-2.tex) external document and shortcuts
\importpaperaux{cpsv16-}{1606.00608/MPDO-22-12-17-2-tnlean}
\newcommand{\cpsvref}[1]{\paperref{cpsv16}{#1}}
\newcommand{\cpsveqref}[1]{\papereqref{cpsv16}{#1}}

% Verdict marker: \gapnote{<kind>}{<status>}. Kind is the policy
% classification (clarification, local-correction, scope-restriction,
% unfaithful, false-source, open-gap); status is open, wip (elimination
% underway), resolved, or historical. Machine-read by the site index;
% prints a verdict line.
\newcommand{\gapnote}[2]{\par\noindent\textbf{Verdict:} #1 (#2).\par}

\importpaperaux{fbc25-}{2502.20257/main-tnlean}

\title{Endpoint-Coordinate Scope of Truncated-Symmetry Unitarity}
\author{}
\date{2026-09-06}

\begin{document}
\maketitle
\gapnote{scope-restriction}{open}

\section{The assertion and its ambient setting}

The defect construction of Ref.~\cite{FrancoRubio2025MPUGauging} starts from a
finite group represented by simple injective matrix product unitaries. In that
setting, a finite-order representing operator has zero source index, and the
paper identifies the two source dimensions with the physical dimension
\cite[line~1547]{FrancoRubio2025MPUGauging}:
\begin{align}
    d_r &= d, & d_l &= d.
    \label{eq:fbc25_truncated_symmetry_source_rank_identification}
\end{align}
The truncated-symmetry display \papereqref{fbc25}{eq:truncsym} then defines
\(U_g^L\) by closing a finite MPU string with the source triangles \(Y_2^g\)
and \(Y_1^g\). Immediately after the display, the source says that \(U_g^L\)
is unitary and acts on local operators in the interior as the global symmetry
does \cite[lines~2062--2099]{FrancoRubio2025MPUGauging}.

Thus the source paragraph has two layers. First, the endpoint contraction is a
unitary. Second, in the finite-group defect setting, this contraction is a
spatially truncated physical symmetry after the endpoint source legs have the
physical dimensions in
(\ref{eq:fbc25_truncated_symmetry_source_rank_identification}).

\section{The endpoint-coordinate theorem}

The formal theorem proves the first layer for a single simple canonical-form-II
MPU tensor.\footnote{Formal declaration:
    \leanid{MPOTensor.IsMPUCanonicalFormII.truncatedSymmetry_isUnitaryBetween}
    in \doclink{TNLean/MPS/MPU/TruncatedSymmetryUnitarity}.}
For \(N\geq0\), the endpoint contraction is the matrix
\begin{align}
    U^N :
    \C^d\otimes(\C^d)^{\otimes N}\otimes\C^d
      &\longrightarrow
    \C^{d_l}\otimes(\C^d)^{\otimes N}\otimes\C^{d_r},
    \label{eq:fbc25_truncated_symmetry_formal_spaces}
\end{align}
with entries
\begin{align}
    (U^N)_{(l,\sigma,r),(a,\tau,b)}
      &=
      \sum_{\alpha,\beta}
      (Y_2)_{l,(a,\alpha)}
      \mathcal U^{\sigma,\tau}_{\alpha,\beta}
      (Y_1)_{r,(\beta,b)}.
    \label{eq:fbc25_truncated_symmetry_formal_entries}
\end{align}
Here \(\sigma,\tau\) are words of length \(N\), and the superscript \(N\)
counts only the bulk tensor letters, not the two endpoint triangles. The
formal conclusion is unitarity between the two coordinate spaces:
\begin{align}
    (U^N)^\dagger U^N
      &= I_{\C^d\otimes(\C^d)^{\otimes N}\otimes\C^d},
    \label{eq:fbc25_truncated_symmetry_formal_isometry} \\
    U^N(U^N)^\dagger
      &= I_{\C^{d_l}\otimes(\C^d)^{\otimes N}\otimes\C^{d_r}}.
    \label{eq:fbc25_truncated_symmetry_formal_coisometry}
\end{align}

The theorem assumes canonical form II and simplicity. The source-gate
equivalence theorem for simple MPUs gives
\cite[lines~532--543 and 563--601]{Cirac2017MPU}
\begin{align}
    d_rd_l &= d^2,
    \label{eq:fbc25_truncated_symmetry_rank_product} \\
    u^\dagger u &= uu^\dagger = I.
    \label{eq:fbc25_truncated_symmetry_source_gate_unitarity}
\end{align}
Together with the movement-gate unitarity of
\papereqref{fbc25}{eq:wLR} \cite[lines~811--869 and 918--1012]{FrancoRubio2025MPUGauging},
these inputs prove
(\ref{eq:fbc25_truncated_symmetry_formal_isometry}) and
(\ref{eq:fbc25_truncated_symmetry_formal_coisometry}) by induction on \(N\).
The rank product in
(\ref{eq:fbc25_truncated_symmetry_rank_product}) makes the two sides of
(\ref{eq:fbc25_truncated_symmetry_formal_spaces}) have equal total dimension:
\begin{align}
    d^{N+2} &= d_l d^N d_r.
    \label{eq:fbc25_truncated_symmetry_total_dimension}
\end{align}
It does not by itself provide the individual endpoint identifications
\begin{align}
    \C^{d_l} &\simeq \C^d, & \C^{d_r} &\simeq \C^d.
    \label{eq:fbc25_truncated_symmetry_missing_endpoint_equivalences}
\end{align}

\section{Growth and what remains outside the theorem}

The two-end source-factorization identity represented by
\papereqref{fbc25}{eq:move_trunc_sym} is already formalized separately. It
adds one bulk site at each end:
\begin{align}
    (U^{N+2})_{(l,i\sigma j,r),(a,b\tau c,e)}
      &=
      \sum_{t,s}
      (w_L)_{(l,i),(a,t)}
      (U^N)_{(t,\sigma,s),(b,\tau,c)}
      (w_R)_{(j,r),(s,e)}.
    \label{eq:fbc25_truncated_symmetry_two_end_growth}
\end{align}
The all-length unitarity proof uses the one-left-end version of this
factorization, together with unitary \(w_L\). This is a proof organization
choice, not a missing two-end factorization.

The statements still outside the endpoint-coordinate theorem are different.
It does not derive the finite-group endpoint dimensions
(\ref{eq:fbc25_truncated_symmetry_source_rank_identification}), choose the
endpoint equivalences in
(\ref{eq:fbc25_truncated_symmetry_missing_endpoint_equivalences}), or prove the
interior action condition stated after \papereqref{fbc25}{eq:truncsym}. The
complementary finite-chain relation is adjacent remaining scope after the
movement display, not part of the literal sentence at lines 2062--2099, and
the endpoint-coordinate theorem does not prove it:
\begin{align}
    U_g^L U_{g^{-1}}^\infty &= U_{g^{-1}}^{\infty-L}.
    \label{eq:fbc25_truncated_symmetry_complementary_relation}
\end{align}
This is \papereqref{fbc25}{eq:complementary} in
Ref.~\cite[lines~2176--2181]{FrancoRubio2025MPUGauging}, and it depends on the
finite-group and gauge choices preceding the defect construction.

\section{Verdict and formal boundary}

The formal result is a complete endpoint-coordinate unitarity theorem: it
proves both Gram identities
(\ref{eq:fbc25_truncated_symmetry_formal_isometry}) and
(\ref{eq:fbc25_truncated_symmetry_formal_coisometry}) for every bulk length.
It is not merely an isometry, and it does not assume supplied unitarity of
the source or movement gates.

It is nevertheless a scope restriction relative to the source paragraph.
The source is discussing a finite-group spatial truncation of a physical
symmetry, whereas the formal theorem stops at unitarity between the input
space in (\ref{eq:fbc25_truncated_symmetry_formal_spaces}) and the endpoint
source-coordinate output space. Recovering the full source paragraph requires
the finite-group source-rank identification, endpoint coordinate choices, the
interior action condition, and the complementary relation
(\ref{eq:fbc25_truncated_symmetry_complementary_relation}). The restricted
theorem may be cited as formalizing only the endpoint-coordinate unitarity
part of the source assertion.

\bibliographystyle{alpha}
\bibliography{references}

\end{document}
